\section*{Hoofdstuk \ref{ch:b3:exploration-seismology} --- Exploratieseismologie}

\begin{solution}{exo:b3:exploration-seismology:1}
$Z = \rho v$: krijt $2000\times2500 = \qty{5.00e6}{kg.m^{-2}.s^{-1}}$, schalie
$2300\times2300 = \qty{5.29e6}{kg.m^{-2}.s^{-1}}$, kalksteen $2600\times4500 =
\qty{11.70e6}{kg.m^{-2}.s^{-1}}$. Krijt--schalie: $R = (5.29 - 5.00)/(5.29 + 5.00)
= 0.29/10.29 = 0.0282$; schalie--kalksteen: $R = (11.70 - 5.29)/(11.70 + 5.29) =
6.41/16.99 = 0.377$. De gereflecteerde energie is $R^2 = 0.142$: de top van de
kalksteen kaatst \qty{14.2}{\percent} van de invallende energie terug, het krijt--schalie
grensvlak minder dan \qty{0.1}{\percent}.
\end{solution}

\begin{solution}{exo:b3:exploration-seismology:2}
$t_0 = 2\times800/2000 = \qty{0.800}{s}$; $t = \sqrt{0.800^2 + (1000/2000)^2} =
\sqrt{0.89} = \qty{0.9434}{s}$; de \omterm{thm:b3:exploration-seismology:nmo}{normal moveout} is \qty{0.1434}{s}. De
benadering geeft $1000^2/(2\times2000^2\times0.800) = \qty{0.1563}{s}$, te
groot met \qty{9}{\percent}: bij een offset groter dan de diepte geldt $x \ll vt_0$ niet
meer.
\end{solution}

\begin{solution}{exo:b3:exploration-seismology:3}
De gemeenschappelijke middelpunten verschuiven over de helft van het ontvangerinterval: het CMP-interval is
$12.5/2 = \qty{6.25}{m}$. De bedekkingsgraad is $240\times12.5/(2\times25) = 60$. Een
bedekkingsgraad van 120 vereist een schot om de $240\times12.5/(2\times120) = \qty{12.5}{m}$.
\end{solution}

\begin{solution}{exo:b3:exploration-seismology:4}
$\lambda = v/f$: \qty{100}{m} in het zandsteen, $\lambda/4 = \qty{25}{m}$;
\qty{150}{m} in de kalksteen, $\lambda/4 = \qty{37.5}{m}$.
Met $\lambda = 2800/30 = \qty{93.3}{m}$ en $z = \qty{1000}{m}$,
$r_F \approx \sqrt{\lambda z/2} = \sqrt{46\,667} = \qty{216}{m}$ (\qty{217}{m}
met de term $\lambda^2/16$).
\end{solution}

\begin{solution}{exo:b3:exploration-seismology:5}
\omterm{thm:b3:exploration-seismology:dix}{Dix-vergelijking}: $v^2 = (2800^2\times0.700 - 2400^2\times0.500)/0.200 = (5\,488\,000 -
2\,880\,000)/0.200 = \qty{1.304e7}{m^2.s^{-2}}$, dus $v = \qty{3611}{m.s^{-1}}$.
De laag wordt in \qty{0.100}{s} in één richting doorlopen: $h = 3611\times0.100 =
\qty{361}{m}$.
\end{solution}

\begin{solution}{exo:b3:exploration-seismology:6}
(a) Laat $d$ de afstand zijn van het middelpunt M tot de reflector, zodat
$t_0 = 2d/v$. De gereflecteerde straal van de bron S naar de ontvanger G is even lang
als de rechte lijn van het beeld S$'$ van S in de reflector naar G. Plaats M in
de oorsprong, de lijn langs de helling; S$'$ ligt op de afstand $2d_S$ van S langs
de normaal, met $d_S = d - (x/2)\sin\varphi$. Het kwadraat van S$'$G is
$x^2 + 4d_S^2 + 4d_Sx\sin\varphi = x^2 + 4d^2 - x^2\sin^2\varphi$, waarbij alle termen
in $dx$ wegvallen, zodat $v^2t^2 = 4d^2 + x^2\cos^2\varphi$ en
$t^2 = t_0^2 + x^2\cos^2\varphi/v^2$. (b) De gather volgt een hyperbool met snelheid
$V = v/\cos\varphi = 2000/\cos\ang{12} = \qty{2045}{m.s^{-1}}$: een helling maakt de
\omterm{def:b3:exploration-seismology:velocities}{stapelsnelheid} te hoog, hier met \qty{2.2}{\percent}.
\end{solution}

\begin{solution}{exo:b3:exploration-seismology:7}
Met $\sqrt{60} = 7.75$. Verdubbeling van de signaal-ruisverhouding vereist vier keer zoveel
sporen: een bedekkingsgraad van 240.
\end{solution}

\begin{solution}{exo:b3:exploration-seismology:8}
De uitgezette diepte $vt/2$ verandert met $\qty{1.5}{km.s^{-1}}\times\qty{0.40}{s}
= \qty{0.60}{km}$ per kilometer: $\tan\theta_a = 0.60$, $\theta_a = \ang{31.0}$.
Volgens de zero-offsetstelling $\sin\theta = 0.60$: $\theta = \ang{36.9}$.
\end{solution}

\begin{solution}{exo:b3:exploration-seismology:9}
Elke piek plaatst een kopie van $w$ gecentreerd erop: $0.3\,w$ op de monsters 0--2 geeft
$(-0.15, 0.3, -0.15)$, en $-0.2\,w$ op de monsters 2--4 geeft $(0.1, -0.2, 0.1)$.
Hun som is $(0, -0.15, 0.3, -0.05, -0.2, 0.1, 0)$, geschreven van één monster
vóór de reeks tot één erna. De reflecties pieken bij de monsters 1 en 3, met
hun ware tekens; ze overlappen op monster 2, waar $-0.15 + 0.1 = -0.05$: een
kleine interferentie, maar twee duidelijk onderscheiden pieken.
\end{solution}

\begin{solution}{exo:b3:exploration-seismology:10}
$\arcsin(1800/2900) = \ang{38.37}$ en $\arcsin(1800/4400) = \ang{24.15}$, dus
$i_c = (38.37 + 24.15)/2 = \ang{31.26}$ en $\theta = (38.37 - 24.15)/2 =
\ang{7.11}$; $v_2 = 1800/\sin\ang{31.26} = \qty{3469}{m.s^{-1}}$. De harmonische
gemiddelde $2\times2900\times4400/(2900 + 4400) = \qty{3496}{m.s^{-1}}$ is
\qty{0.8}{\percent} hoger; vermenigvuldigd met $\cos\theta = 0.9923$ geeft het
\qty{3469}{m.s^{-1}}, de gebruikelijke benadering voor kleine hellingen
$v_2 \approx \cos\theta\,2v_dv_u/(v_d + v_u)$.
\end{solution}

\begin{solution}{exo:b3:exploration-seismology:11}
(a) Het samengestelde signaal is $R[w(t) - w(t - \Delta t)] \approx R\,\Delta t\,w'(t)$ door
de eerste-orde Taylorontwikkeling. (b) De amplitude is evenredig met
$\Delta t$, dus een laag die dunner is dan de \omterm{def:b3:exploration-seismology:resolution}{tuningdikte} geeft nog steeds een signaal,
met de vorm van $w'$, maar haar top en basis kunnen niet als twee events worden aangewezen: ze
wordt gedetecteerd, niet opgelost. Als $R$ en de \omterm{def:b3:exploration-seismology:wavelet}{seismische wavelet} bekend zijn, bijvoorbeeld uit een
boorput, geeft de amplitude $\Delta t$ en dus de dikte. (c) De amplitude is
evenredig met de dikte: de verhouding is $5/10 = 1/2$.
\end{solution}

\begin{solution}{exo:b3:exploration-seismology:12}
(a) $t(x) = (2/v)\sqrt{z^2 + x^2}$: $t(0) = 2\times600/3000 = \qty{0.400}{s}$ en
$t(300) = (2/3000)\sqrt{600^2 + 300^2} = \qty{0.4472}{s}$. (b) $t^2 = t_0^2 +
x^2/(v/2)^2$, de NMO-vorm met snelheid $v/2$. Een reflectie met dezelfde
toptijd, op een zero-offsetsectie, zou helemaal geen moveout vertonen; vergeleken met een
NMO-hyperbool met dezelfde $t_0$ in een gather, is de moveout van de diffractie
$x^2/(2(v/2)^2t_0)$ is vier keer zo groot, omdat de tijd tweemaal van de afstand
tot de verstrooier afhangt, heen en terug. (c) Migratie telt de
sectie langs deze hyperbool op en legt de energie in de top: de diffractie
klapt samen tot een punt bij $x = 0$, $z = \qty{600}{m}$.
\end{solution}

\begin{solution}{pb:b3:exploration-seismology:1}
\textbf{1.} Vanuit W, $400/(486.1 - 425.2) = 400/0.0609 = \qty{6568}{m.s^{-1}}$;
vanuit O, $400/(486.1 - 372.4) = 400/0.1137 = \qty{3518}{m.s^{-1}}$. De
kopgolf vanuit W lijkt snel: W is het eind van de helling naar beneden, en zijn schoten lopen de helling op; de
refractor stijgt van W naar O.

\textbf{2.} $\arcsin(2400/3518) = \ang{43.02}$, $\arcsin(2400/6568) =
\ang{21.43}$: $i_c = \ang{32.22}$, $\theta = \ang{10.79}$.

\textbf{3.} $v_2 = 2400/\sin\ang{32.22} = \qty{4501}{m.s^{-1}}$: binnen het
bereik van kalksteen, de top van de Jura-kalksteen.

\textbf{4.} De intercepts zijn $t - x/v_{\text{app}}$ op \qty{1200}{m}:
$t_{i,W} = 0.4861 - 1200/6568 = \qty{0.3034}{s}$ en $t_{i,E} = 0.4861 -
1200/3518 = \qty{0.1450}{s}$. Met $\cos i_c = 0.8460$, $h = v_1t_i/(2\cos
i_c)$: $h_W = 2400\times0.3034/1.692 = \qty{430.4}{m}$ en $h_E =
\qty{205.7}{m}$.

\textbf{5.} Delen door $\cos\theta = 0.9823$: \qty{438.1}{m} onder W en
\qty{209.4}{m} onder O. Dan is $(438.1 - 209.4)/1200 = 0.1906 =
\tan\ang{10.79}$, de helling uit vraag 2.

\textbf{6.} De schalie levert geen kopgolf op: ze is onzichtbaar, en het
onderzoek telt haar mee bij het krijt. Hier is $v_1 = \qty{2.40}{km.s^{-1}}$ gegeven als
het gemiddelde van de twee; had de directe-golfsnelheid van het krijt, \qty{2.50}{km.s^{-1}},
voor de hele bedekking gebruikt, dan zouden alle dieptes te groot uitvallen.

\textbf{7.} $V = 600/\sqrt{0.5437^2 - 0.4878^2} = 600/0.2401 =
\qty{2499}{m.s^{-1}}$ en $V = 600/\sqrt{0.6240^2 - 0.5748^2} =
\qty{2471}{m.s^{-1}}$. De tweede middelt het krijt met de langzamere schalie
eronder: het gemiddelde daalt.

\textbf{8.} \qty{2811}{m.s^{-1}} en \qty{2829}{m.s^{-1}}.

\textbf{9.} Laag 1: \qty{2499}{m.s^{-1}}. Laag 2: $v^2 = (2471^2\times0.5748
- 2499^2\times0.4878)/0.0870$, $v = \qty{2308}{m.s^{-1}}$. Laag 3:
$v^2 = (2811^2\times0.6548 - 2471^2\times0.5748)/0.0800$,
$v = \qty{4561}{m.s^{-1}}$. Laag 4: $v^2 = (2829^2\times0.7348 -
2811^2\times0.6548)/0.0800$, $v = \qty{2972}{m.s^{-1}}$.

\textbf{10.} Krijt (\qty{2499}{m.s^{-1}}), schalie (\qty{2308}{m.s^{-1}}),
kalksteen (\qty{4561}{m.s^{-1}}), zandsteen (\qty{2972}{m.s^{-1}}), elk binnen
het gepubliceerde bereik. De reflecties komen van het grensvlak krijt--schalie,
de top van de kalksteen, het grensvlak kalksteen--zandsteen en de
discordantie op de top van het schistbasement.

\textbf{11.} Krijt en schalie hebben vrijwel dezelfde impedantie: $R = 0.028$
(oefening 1), tegenover 0.38 voor de top van de kalksteen.

\textbf{12.} De hyperbool met snelheid $V_{\text{rms}}$ geldt alleen tot de
tweede orde in de offset; bij offsets van \qty{2}{km}, twee keer de diepte
van de diepste reflector, wijken de tijden ervan af, en een hyperbool die op
het hele spreidingsgebied wordt gefit, zou verkeerde snelheden geven, en Dix-intervallen nog slechter. De verre offsets
rekken bovendien de wavelets op bij de correctie.

\textbf{13.} Heen-en-terugtijden in de lagen: \qty{0.4878}{s}, \qty{0.0870}{s},
\qty{0.0800}{s}, \qty{0.0800}{s}. Diktes $v\,\Delta t/2$: \qty{609.5}{m},
\qty{100.4}{m}, \qty{182.4}{m}, \qty{118.9}{m}.

\textbf{14.} Dieptes onder het maaiveld: \qty{609.5}{m}, \qty{709.9}{m},
\qty{892.3}{m}, \qty{1011.2}{m}. De schalie, kalksteen en zandsteen komen
binnen \qty{0.4}{m}, \qty{2.4}{m} en \qty{1.1}{m} van hun diktes uit, waarbij
de fout groeit in de snelle, dunne kalksteen, waar een kleine tijdsfout het
meest kost.

\textbf{15.} $\bar v = 2\times1011.2/0.7348 = \qty{2752}{m.s^{-1}}$. De
\omterm{def:b3:exploration-seismology:velocities}{RMS-snelheid} weegt elke laag met $v^2$ en is nooit kleiner dan het
tijdgewogen gemiddelde: $2829\times0.7348/2 = \qty{1039}{m}$, \qty{28}{m} te diep.

\textbf{16.} $\lambda = 2308/30 = \qty{76.9}{m}$, $\lambda/4 = \qty{19.2}{m}$;
de Ricker-schatting is $2308/(5.2\times30) = \qty{14.8}{m}$. De schalie is
\qty{100}{m} dik, vijf keer haar \omterm{def:b3:exploration-seismology:resolution}{tuningdikte}: haar top en basis worden
opgelost. Ze zou volgens de kwart-golflengteregel afstemmen bij $f = v/(4b) =
2308/400 = \qty{5.77}{Hz}$.

\textbf{17.} $\lambda = 2829/30 = \qty{94.3}{m}$, $z = \qty{1011}{m}$:
$r_F \approx \sqrt{94.3\times1011/2} = \qty{218}{m}$ (\qty{220}{m} met de
term $\lambda^2/16$).

\textbf{18.} $t(x)^2 = t_0^2 + (2x/v)^2$, dus $v = 2x/\sqrt{t^2 - t_0^2} =
400/\sqrt{0.5732^2 - 0.550^2} = 400/0.1614 = \qty{2478}{m.s^{-1}}$.

\textbf{19.} $z = vt_0/2 = 2478\times0.550/2 = \qty{681}{m}$.

\textbf{20.} De uitgezette diepte $vt/2$ verandert met $1.24\times0.167 = 0.207$
per eenheid lengte: $\theta_a = \arctan0.207 = \ang{11.70}$. De ware helling
heeft $\sin\theta = 0.207$: $\theta = \ang{11.95}$, de helling van de flank.

\textbf{21.} $d = 2480\times0.655/2 = \qty{812}{m}$; migratie verplaatst de
reflectie $d\sin\theta = 812\times0.207 = \qty{168}{m}$ de helling op, richting de
kam van de antiklinaal, en omhoog van \qty{812}{m} naar $d\cos\theta =
\qty{795}{m}$.

\textbf{22.} $2500\times0.160/2 = \qty{200}{m}$.

\textbf{23.} Uit vragen 13 en 14 ligt de basisdiscordantie
\textbf{\qty[text-series-to-math=true]{1011}{m}} onder het maaiveld bij CMP S.
\end{solution}
